% !TeX program = xelatex
% Overleaf Compiler: XeLaTeX
% 注意：本模板中的姓名、学校、公司、经历、成绩和联系方式均为虚构示例，仅用于展示排版。
\documentclass[9pt,a4paper]{ctexart}
\usepackage[left=15mm,right=15mm,top=9mm,bottom=9mm]{geometry}
\usepackage{tikz,xcolor,graphicx,array,tabularx,enumitem}
\pagestyle{empty}
\setlength{\parindent}{0pt}
\setlength{\parskip}{0pt}
\linespread{1.025}
\setlength{\tabcolsep}{0pt}

\definecolor{resbg}{HTML}{F7F8F8}
\definecolor{resline}{HTML}{DADFE1}
\definecolor{restext}{HTML}{222222}
\definecolor{resarrow}{HTML}{FFFDF5}
\definecolor{secinfo}{HTML}{DEA665}
\definecolor{secedu}{HTML}{7FAFCB}
\definecolor{secproj}{HTML}{4AAAC3}
\definecolor{seccampus}{HTML}{69B79C}
\definecolor{secaward}{HTML}{9DBD64}
\definecolor{secself}{HTML}{B18FC5}

\newcommand{\ResumeTitle}{\noindent{\bfseries\fontsize{12.5}{15}\selectfont 个人简历}\par\vspace{1.7mm}}
\newcommand{\ResumeBanner}[2]{%
  \begin{tikzpicture}[baseline=(current bounding box.north)]
    \fill[resbg,rounded corners=1.2mm] (0,0) rectangle (5.15,0.98);
    \fill[#1,rounded corners=0.35mm] (0.07,0.04) rectangle (0.43,0.94);
    \node[text=resarrow,font=\bfseries\fontsize{9}{10}\selectfont,inner sep=0pt] at (0.25,0.49) {\textgreater};
    \node[anchor=west,text=restext,font=\bfseries\fontsize{13}{15}\selectfont,inner sep=0pt] at (0.72,0.49) {#2};
    \draw[resline,line width=0.35pt] (0.48,0.04)--(5.08,0.04);
  \end{tikzpicture}%
}
\newcommand{\ResumeSection}[2]{\par\vspace{1.7mm}\noindent\ResumeBanner{#1}{#2}\par\vspace{1.0mm}}
% 使用通用插画头像，避免误放真实个人照片
\newcommand{\ResumePhoto}{\includegraphics[width=26mm,height=34mm]{generic-photo.png}}
\newcommand{\ResumeBasicBlock}[2]{%
  \noindent\begin{minipage}[t]{0.79\linewidth}\vspace{0pt}\ResumeBanner{secinfo}{基本信息}\par\vspace{2mm}#1\end{minipage}\hfill%
  \begin{minipage}[t]{0.19\linewidth}\vspace{0pt}\raggedleft #2\end{minipage}\par
}
\newcommand{\Edu}[3]{%
  \noindent\begin{tabularx}{\linewidth}{@{}>{\bfseries\raggedright\arraybackslash}p{0.19\linewidth}>{\raggedright\arraybackslash}X>{\bfseries\raggedleft\arraybackslash}p{0.22\linewidth}@{}}
  #1 & \hspace{2mm}#2 & #3
  \end{tabularx}\par\vspace{1.0mm}%
}
\newcommand{\WorkExperience}[4]{%
  \noindent\begin{tabularx}{\linewidth}{@{}>{\bfseries\raggedright\arraybackslash}p{0.42\linewidth}>{\bfseries\centering\arraybackslash}p{0.28\linewidth}>{\bfseries\raggedleft\arraybackslash}X@{}}
  #2 & #3 & #1
  \end{tabularx}\par\vspace{0.45mm}#4\par\vspace{0.65mm}%
}
\newcommand{\ResumeLine}[2]{\noindent\textbf{#1}：#2\par\vspace{0.3mm}}

\begin{document}
\ResumeTitle
\ResumeBasicBlock{%
  姓\quad 名：林知夏\qquad 目标岗位：\textbf{产品经理实习生}\par\vspace{1.8mm}
  出生年月：2002年06月\enspace\textcolor{resline}{|}\enspace 邮箱：lin.zhixia@example.com\par\vspace{1.8mm}
  电话：13800000000\enspace\textcolor{resline}{|}\enspace 微信：linzhixia\_pm
}{\ResumePhoto}

\ResumeSection{secedu}{教育背景}
\Edu{东川大学}{信息管理与信息系统\quad 硕士（GPA：3.78/4.00，专业前15\%）}{2024.09--2027.06}
\Edu{江城理工大学}{信息管理与信息系统\quad 本科（GPA：3.62/4.00，专业前20\%）}{2020.09--2024.06}
\ResumeLine{奖学金}{硕士学业二等奖学金（2024--2025）；硕士学业一等奖学金（2025--2026）；本科校级二等奖学金（2021--2022、2022--2023）。}
\ResumeLine{荣誉竞赛}{校级优秀学生（2022）；“互联网+”校级创新创业大赛二等奖（2023）；优秀学生干部（2023）。}
\ResumeLine{语言证书}{大学英语四级 CET-4（548分）；大学英语六级 CET-6（492分）；普通话二级甲等。}

\ResumeSection{seccampus}{实习与工作经历}
\WorkExperience{2026.03--2026.06}{星河科技有限公司}{产品经理实习生}{%
  \ResumeLine{需求分析}{参与企业协作工具的用户访谈与竞品分析，整理12份访谈记录和6款竞品功能对比，归纳任务提醒、权限配置等高频需求。}
  \ResumeLine{产品设计}{协助撰写产品需求文档（PRD），绘制核心流程图及低保真原型，跟进设计、研发和测试评审，推动3项功能进入迭代。}
  \ResumeLine{数据复盘}{使用SQL和Excel跟踪功能使用率、任务完成率等指标，整理上线后数据周报，为后续交互优化提供依据。}
}
\vspace{1.2mm}
\WorkExperience{2025.07--2025.09}{云帆网络科技有限公司}{产品运营实习生}{%
  \ResumeLine{用户运营}{协助维护新用户引导与产品帮助中心，分析用户反馈并归纳高频问题，推动更新8篇帮助文档。}
  \ResumeLine{活动复盘}{跟踪线上专题活动的曝光、点击和转化数据，使用Excel制作周报与活动复盘表，提出页面文案和入口位置优化建议。}
}

\ResumeSection{secproj}{项目经历}
\WorkExperience{2025.10--2026.01}{校园二手交易平台改版项目}{产品负责人（课程项目）}{%
  \ResumeLine{用户研究}{面向在校学生开展问卷调研并访谈15名潜在用户，识别商品发布流程复杂、交易信任不足等主要问题。}
  \ResumeLine{方案设计}{使用Figma绘制商品发布、搜索筛选和信用评价原型，编写需求文档并拆分功能优先级，组织4人小组完成可交互原型。}
  \ResumeLine{项目结果}{通过5名用户进行可用性测试，依据任务完成情况优化关键页面流程，形成需求文档、原型及测试报告。}
}
\vspace{1.0mm}
\WorkExperience{2025.03--2025.06}{校园活动报名小程序}{产品策划（课程项目）}{%
  \ResumeLine{功能规划}{梳理活动发布、报名审核、消息通知和签到流程，绘制用户流程图及页面原型，完成项目汇报与演示。}
}

\ResumeSection{seccampus}{校园经历}
\WorkExperience{2023.09--2024.06}{江城理工大学学生会}{宣传部副部长}{%
  \ResumeLine{活动组织}{协助策划迎新季及校园文化节宣传工作，协调6名部门成员完成活动海报、推文和现场素材整理。}
  \ResumeLine{内容运营}{参与运营学院微信公众号，规划选题、撰写推文并跟踪阅读数据；学年内参与发布推文18篇，单篇最高阅读量约1,200次。}
}

\ResumeSection{secself}{技能特长}
\ResumeLine{产品工具}{熟悉Axure、Figma进行原型设计，能够使用XMind梳理需求与业务流程，了解PRD撰写和需求优先级管理。}
\ResumeLine{数据分析}{掌握Excel函数、数据透视表和图表制作，了解SQL基础查询，能够进行基础指标统计与活动复盘。}
\ResumeLine{办公与语言}{熟练使用PowerPoint、Word；CET-4、CET-6，具备英文资料阅读与基础书面沟通能力。}
\ResumeLine{综合能力}{具备用户视角、逻辑分析、跨团队沟通和项目协作意识，能够使用ChatGPT等AI工具辅助资料整理与方案梳理。}

\end{document}
